% !TeX root = ../../surface-lagrangian-torus-fibration.tex
% 2026-11-21 / v1.3: 二段階平均，固定残余領域の法線エネルギーと接続上界．
\section{共通の normal circle に適合する二段階平均}
\label{ke:staged-neck-section}

三重点の torus 平均を，neck 全体で共通に用いる normal circle の平均と，
end collar 上の残る circle の平均に分ける．
この順序によって，完成後の end tori が一つの actual Hamiltonian circle を含むことを
厳密に保証できる．単一 primitive を用いると normal moment の加法定数と
全 marked cycle の作用周期も同時に固定される．
平均化と Moser 法そのものの先行研究は
\cite[Section~3]{RuanHMinimal}を参照されたい．
以下の二段階構成と定量的適用の包括的な新規性確認は未完了である．

本節の円作用の周期はすべて $1$ とし，moment の符号は
\[
 \iota_V\omega=-d\mu
\]
に統一する．従って moment は一周の primitive 積分と同じ規格化である．
周期 $2\pi$ の生成場による通常の action 座標を用いる場合は，
生成場・moment・周期積分の規格化を同時に変更する．
また，以下の flow は指定する内側領域から出発して共通の外側領域に留まると仮定する．
作用・形式が一致する開集合を出る trajectory の一致は，局所的な ODE の一意性からは従わない．

\subsection{一つの円の平均と primitive の正確な輸送}

$U$ を複素曲面，$\omega=d\lambda$ を K\"ahler form とし，
$\rho_s$ を $U$ 上の自由な正則円作用とする．
生成場を $V$，平均と chain homotopy を
\[
 \mathcal A\xi=\int_0^1\rho_s^*\xi\,ds,\qquad
 K\xi=\int_0^1(1-s)\rho_s^*(\iota_V\xi)\,ds
\]
と書く．Centered kernel $1/2-s$ に置き換えた場合も，
以下の恒等式と結論は同じである．

\begin{lemma}[円平均の exact Moser transport]
\label{ke:circle-exact-transport}
$\eta=\mathcal A\omega$，
$\omega_s=(1-s)\eta+s\omega$，
$\lambda_s=(1-s)\mathcal A\lambda+s\lambda$ とする．
$\beta=K\omega$ および
\[
 \iota_{Y_s}\omega_s=\beta,\qquad \phi_0=\mathrm{id}
\]
による flow が全時刻で定義される $\rho$-飽和した内側領域を取る．この領域では
\begin{equation}\label{ke:eq:circle-exact-transport}
 \phi_s^*\omega_s=\eta,\qquad
 \phi_s^*\lambda_s=\mathcal A\lambda+dq_s,
 \quad
 q_s=\int_0^s\phi_\tau^*
       \bigl(\lambda_\tau(Y_\tau)-K\lambda\bigr)\,d\tau .
\end{equation}
$\mu=(\mathcal A\lambda)(V)$ は $\eta$ の Hamiltonian moment である．
$\phi=\phi_1$ によって移した円作用
$\widehat\rho_s=\phi\circ\rho_s\circ\phi^{-1}$ は actual $\omega$ に関して Hamiltonian で，
その moment は
\begin{equation}\label{ke:eq:circle-actual-moment}
 \widehat\mu=\mu\circ\phi^{-1}.
\end{equation}
この moment の値は，対応する actual 円軌道上の $\lambda$ の一周積分に等しい．
\end{lemma}
\begin{proof}
Cartan の公式と部分積分により
\[
 dK+Kd=\mathcal A-I.
\]
従って
\[
 d\beta=\eta-\omega,\qquad
 \beta=\mathcal A\lambda-\lambda-d(K\lambda).
\]
正則作用による平均とその凸結合は K\"ahler form なので非退化である．
Moser の微分式から最初の pullback identity を得る．
Primitive については
\[
 \partial_s\lambda_s+\mathcal L_{Y_s}\lambda_s
 =\lambda-\mathcal A\lambda+\beta
     +d\bigl(\lambda_s(Y_s)\bigr)
 =d\bigl(\lambda_s(Y_s)-K\lambda\bigr)
\]
であり，flow で引き戻して積分すれば
\eqref{ke:eq:circle-exact-transport}の第二式となる．

$\mathcal A\lambda$ は円作用に不変なので
\[
 \iota_V\eta=-d\bigl((\mathcal A\lambda)(V)\bigr)=-d\mu.
\]
これを $\phi$ で移すと\eqref{ke:eq:circle-actual-moment}を得る．
周期 $1$ の軌道 $s\mapsto\rho_s(x)$ 上で
$(\mathcal A\lambda)(V)=\mu(x)$ は一定である．
Exact differential の周期は零なので，
$\int_{\phi(\rho_\bullet(x))}\lambda=\mu(x)$ である．
\end{proof}

\begin{lemma}[Actual circle の短さによる双方向 Moser buffer]
\label{ke:short-circle-actual-buffer}
上記の正則円作用と K\"ahler form に対し，作用不変な開集合 $U$ 全体で
$|V|_{g_\omega}\le\varepsilon$ と仮定する．
この補題の評価には $\omega$ の exactness，作用不変な背景計量，
tensor の二側比較は不要である．
$w(a)=1-a$ または $w(a)=1/2-a$ とし，
\[
 \beta=\int_0^1 w(a)\rho_a^*(\iota_V\omega)\,da,
 \quad \eta=\mathcal A\omega,
 \quad \omega_s=(1-s)\eta+s\omega,
 \quad \iota_{Y_s}\omega_s=\beta
\]
と置く．$c_w=(\pi/2)\|w\|_{L^2(0,1)}$ とすれば，
\begin{equation}\label{ke:eq:circle-actual-speed}
 |Y_s|_{g_\omega}
 \le\frac{\varepsilon\|w\|_{L^2(0,1)}}{2\sqrt{s(1-s)}}
 \qquad (0<s<1).
\end{equation}
従って，任意の二つの時刻を結ぶ存在する trajectory の actual 長さは
\begin{equation}\label{ke:eq:circle-actual-length}
 \operatorname{Length}_{g_\omega}(\phi_{s,s_0}(x))
 \le c_w\varepsilon,
 \qquad
 c_w=\begin{cases}
 \pi/(2\sqrt3),&w(a)=1-a,\\
 \pi/(4\sqrt3),&w(a)=1/2-a
 \end{cases}
\end{equation}
である．ここで時間を逆向きに辿る場合も同じ評価を用いる．

$K\subset U$ を compact な $\rho$-飽和集合とする．
$K\subset D\Subset U$ なる開集合 $D$ が存在し，
\begin{equation}\label{ke:eq:circle-actual-buffer}
 \delta:=d_{g_\omega}^{U}(K,\partial D)>c_w\varepsilon
\end{equation}
を満たすとする．$d_{g_\omega}^{U}$ は $U$ 内の経路による actual 距離である．
このとき $K$ のある $\rho$-飽和開近傍からの forward flow は全時刻で存在する．
また，時刻 $1$ に $K$ の点を終点として指定する backward flow も時刻 $0$ まで存在する．
従って最大の time-one flow の定義域と像はともに $K$ を含む．
各円軌道 $\mathcal O\subset K$ に対して
\begin{equation}\label{ke:eq:circle-actual-diameter}
 \operatorname{diam}_{d_{g_\omega}^{U}}\phi_{1,0}(\mathcal O)
 \le (2c_w+1/2)\varepsilon.
\end{equation}
ここで直径は ambient 距離による．
像の円自身の intrinsic diameter，flow の $C^1$ 評価，
元の neck base projection の保存は結論に含めない．
\end{lemma}
\begin{proof}
$g_a=\rho_a^*g_\omega$ と置く．正則性により
$g_\eta=\int_0^1g_a\,da$ であり，$d\rho_a(V)=V$ と仮定から
$g_a(V,V)\le\varepsilon^2$ である．任意の接ベクトル $X$ に対し
\begin{align*}
 |\beta(X)|
 &\le\varepsilon\int_0^1 |w(a)|\,|X|_{g_a}\,da\\
 &\le\varepsilon\|w\|_{L^2(0,1)}
       \left(\int_0^1g_a(X,X)\,da\right)^{1/2}
 =b|X|_{g_\eta},\qquad b=\varepsilon\|w\|_{L^2(0,1)}.
\end{align*}
共通複素構造を $J$ とする．Moser 方程式から
$g_{\omega_s}(Y_s,Y_s)=\beta(JY_s)$ なので，
$r=|Y_s|_{g_\eta}$，$q=|Y_s|_{g_\omega}$ と置けば
\[
 (1-s)r^2+s q^2\le br.
\]
平方完成によって
\[
 s q^2\le br-(1-s)r^2\le\frac{b^2}{4(1-s)}
\]
となり，\eqref{ke:eq:circle-actual-speed}を得る．
$\int_0^1(s(1-s))^{-1/2}\,ds=\pi$ を用いて
\eqref{ke:eq:circle-actual-length}が従う．
端点で表示した上界は発散するが可積分である．
$Y_s$ 自身は K\"ahler 形式 $\omega_s$ の非退化性により
$[0,1]\times U$ 上で滑らかに定義されている．

任意の時刻に $K$ から出発する trajectory が初めて $D$ の境界へ達すれば，
そこまでの長さは $\delta$ 以上になる．
これは\eqref{ke:eq:circle-actual-length}に反する．
従って trajectory は $\overline D\Subset U$ に留まり，
compact 集合上の滑らかな時間依存 ODE の延長定理によって
順方向・逆方向とも $[0,1]$ 全体へ延びる．
Strict な余裕と $K$ の compact 性により，この結論は $K$ の開近傍にも成立する．
$K$ と作用群が compact なので，その近傍に含まれる
$\rho$-飽和開近傍を選べる．Backward flow の時刻 $0$ の点は forward flow の定義域に属するから，
最大 flow の像も $K$ を含む．

元の円軌道の任意の二点は高々半周で結べ，その長さは $\varepsilon/2$ 以下である．
両端にそれぞれ長さ $c_w\varepsilon$ 以下の Moser trajectory を付ければ
\eqref{ke:eq:circle-actual-diameter}となる．
\end{proof}

\begin{corollary}[同じ compact 領域を覆う actual 円作用]
\label{ke:short-circle-image-buffer}
補題\ref{ke:short-circle-actual-buffer}の仮定に加えて
\begin{equation}\label{ke:eq:circle-image-buffer}
 \delta>(2c_w+1/2)\varepsilon
\end{equation}
とする．このとき time-one Moser map $\phi$ の定義域には
$\rho$-飽和した開集合 $O\supset K$ を取ることができ，その像も $K$ を含む：
\[
 K\subset O,\qquad K\subset\phi(O).
\]
従って $\widehat\rho=\phi\rho\phi^{-1}$ は $K$ の開近傍上の円作用となる．
さらに $\omega=d\lambda$ の場合は補題\ref{ke:circle-exact-transport}の
actual Hamiltonian action と moment の結論がその近傍で成立する．
\end{corollary}
\begin{proof}
$F$ を，時刻 $0$ からの trajectory が全時刻で $D$ 内に留まる点の集合とする．
これは開集合であり，time-one map は $F$ 上の滑らかな埋込である．
その最大の $\rho$-飽和部分
\[
 O=\{y\in F:\rho_a(y)\in F\text{ for every }a\in\mathbb R/\mathbb Z\}
\]
も開である．実際，ある軌道が開集合 $F$ に含まれれば，軌道の compact 性と
作用の連続性によりその近くの全軌道も $F$ に含まれる．
補題の forward 結論と $K$ の飽和性から $K\subset O$ である．

$x\in K$ を取り，backward flow による時刻 $0$ の点を $y$ とする．
$d_{g_\omega}^{U}(x,y)\le c_w\varepsilon$ である．
任意の $z=\rho_a(y)$ は $y$ から高々半周で結べるので
\[
 d_{g_\omega}^{U}(x,z)\le(c_w+1/2)\varepsilon.
\]
$y$ から $z$ への短い軌道弧も $D$ を出られない．
もし出れば，$x$ から $\partial D$ までの経路の長さが
$(c_w+1/2)\varepsilon<\delta$ となるからである．
$z$ からの forward Moser trajectory が $D$ を出るなら，
先ほどの $x$ から $z$ までの経路と合わせて
$x$ から $\partial D$ まで長さ $(2c_w+1/2)\varepsilon$ 以下の経路ができる．
これは\eqref{ke:eq:circle-image-buffer}に反する．
従って全ての $z=\rho_a(y)$ は $F$ に属し，$y\in O$ となる．
$x=\phi(y)$ だから $K\subset\phi(O)$ が従う．
最後の主張は既存の exact transport identity を適用すればよい．
\end{proof}

\begin{remark}[Neck への適用範囲]
上記の二つの結果は，actual normal generator の一様な短さと
actual 距離の buffer から normal averaging を実行する条件である．
Rank-one 評価と volume 比較から短さを得る場合も，
source-free neck 全体でこれらの入力が一様であることを別に証明する必要がある．
複数の triple end における同一作用，reduced leaves，markings の同時整合，
source・branch・core への接着や全曲面の GH 収束は，この補題の結論ではない．
\end{remark}

\subsection{固定残余領域上の法線作用の平均エネルギー}

\begin{proposition}[法線方向の積分評価と短い円軌道]
\label{ke:neck-integrated-normal-energy}
Source と三重点から離れた固定座標で $uv=t$，$z\in D$ とし，
$|u|,|v|<r$ に対して $s=-\log(|u|^2/r^2)\in(0,M_t)$ と置く．
ここで $M_t=-\log(|t|^2/r^4)$ である．
固定した $Q\Subset D$ 上で，actual negative KE form $\omega_t$ の体積が
annulus--disc complete KE 積計量の体積と $a_0\le s\le M_t-a_0$ で
一様に二側比較可能とする（$a_0>0$ は固定）．
$V=2\pi\partial_{\arg u}$，$\rho$ をその周期 $1$ の作用とし，
\[
 E_t(s,z)=\int_0^1 |V|_{g_t}^2(\rho_a(u,z))\,da
\]
と置く．$2a_0\le m\le M_t/2$ ならば，全 $z\in Q$ について
\begin{equation}\label{ke:eq:normal-energy-mass}
 \int_m^{M_t-m} E_t(s,z)\,ds\le C/m .
\end{equation}
定数 $C$ は $t,m,z$ に依存しない．任意の $A>0$ について
$[m,M_t-m]$ 内で $E_t>A/m^2$ となる深度の Lebesgue 測度は
$Cm/A$ 以下である．それ以外の深度では actual 円軌道の長さは
$\sqrt A/m$ 以下である．
\end{proposition}

\begin{proof}
$\Omega=du/u\wedge dz$，$F_t=\log(\omega_t^2/|\Omega|^2)$ とし，
$\bar F_t$ を $u$ の角度について平均する．$|\Omega|^2$ は
$\Omega$ の定める正の体積形式で，その定数規格化は以下に影響しない．
$d^c=\tfrac{\sqrt{-1}}2(\bar\partial-\partial)$ の規格化では
$dd^cF_t=\omega_t$ である．$d(s)=\min(s,M_t-s)$ とすると，
積計量の体積密度は $z\in Q$ 上で
$(\pi/M_t)^2\csc^2(\pi s/M_t)\asymp d(s)^{-2}$ に比較可能だから
\[
 \bigl|\bar F_t(s,z)+2\log d(s)\bigr|\le C_0 .
\]
局所対数座標 $\xi=\log u$ で $s=-2\Re\xi+2\log r$ である．
$\mathcal A\omega_t=dd^c\bar F_t$ の $\xi\bar\xi$ 係数から
$E_t=8\pi^2\partial_s^2\bar F_t$ を得る．特に $\bar F_t$ は $s$ の凸関数である．

$s\in[m,M_t-m]$ で $s\pm m/2$ を取ると，三点の $d$ の値の比は
$1/2$ と $3/2$ の間にある（$d(s)$ を分母とする）．
従って前進・後退の差商は絶対値 $C_1/m$ 以下であり，
凸性から $|\partial_s\bar F_t|\le C_1/m$ を得る．
両端における導関数の差を取れば
\eqref{ke:eq:normal-energy-mass} が従う．
測度評価は Markov の不等式，長さ評価は
$\int_0^1 |V|_{g_t}(\rho_a x)\,da\le E_t(s,z)^{1/2}$ による．
\end{proof}

上の体積仮定は，一般位置全平面退化の固定 source-free pair 座標では，
関連 KE 主稿の \texttt{neck:prop:volume-envelope} から従う．
例えば $m=R_t$，$A=R_t$ とすれば，各固定残余点で中央 strip のうち
深度の測度 $O(1)$ の例外を除いて円軌道長は $O(R_t^{-1/2})$ である．
この例外集合は残余点と $t$ に依存する．

\begin{corollary}[鋭い係数を持つ同一残余点での法線接続上界]
\label{ke:neck-axial-upper}
命題\ref{ke:neck-integrated-normal-energy}の座標で $z\in Q$ を固定する．
$2a_0\le a\le b\le M_t-2a_0$ とする．ある一定の角度で深度 $a$ から $b$ へ
進む動径弧の actual 長さは
\[
 \int_a^b \frac{ds}{d(s)}+C
\]
以下である．従って両端円の任意の点を，同じ $z$ を保つ経路で
この上界に $\tfrac12(\sqrt{E_t(a,z)}+\sqrt{E_t(b,z)})$ を加えた
長さ以下で結べる．
\end{corollary}

\begin{proof}
まず $b\le M_t/2$ とする．$\bar F_t=-2\log s+O(1)$ で，
命題の差商評価に各点で $m=d(s)=s$ を用いると $|s\partial_s\bar F_t|\le C_1$
である．部分積分して
\[
 \int_a^b s\partial_s^2\bar F_t\,ds
 =[s\partial_s\bar F_t]_a^b-[\bar F_t]_a^b
 \le 2\log(b/a)+C_2 .
\]
$JV=4\pi\partial_s$ と $J$ 不変性から，
$g_t(\partial_s,\partial_s)$ の角度平均は
$E_t/(16\pi^2)=\tfrac12\partial_s^2\bar F_t$ である．
従って Cauchy--Schwarz により動径弧長の角度平均は
\[
 \int_a^b \sqrt{\tfrac12\partial_s^2\bar F_t}\,ds
 \le \sqrt{\log(b/a)\bigl(\log(b/a)+C_2/2\bigr)}
 \le \log(b/a)+C_2/4 .
\]
上側半区間では $s$ を $M_t-s$ に置き換える．
中点を横切る場合は角度平均を取った後で両半区間の上界を加えるので，
全区間に共通する一つの角度を選べる．
両端では円の短い方の弧を付加する．
\end{proof}

この接続上界は正規化前の距離に対し加法誤差 $O(1)$ を持つ．
両端の円エネルギーも有界な場合は，正規化した法線接続費用が
$R_t^{-1}\int_a^b ds/d(s)+O(R_t^{-1})$ 以下となる．
この評価だけでは異なる残余点間の portal 接続費用や，全曲面を通る経路の下界は得られない．

\begin{proposition}[固定残余基底上の moving deep moment strip]
\label{ke:neck-deep-moment-strip}
関連 KE 主稿の \texttt{pair:prop:bundle-conic} の固定 source-free compact
基底 $Q\Subset D^\circ$ 上で，双対 Hermitian norms を用いて
$uv=t$，$s=-\log(\|v\|^2/r^2)$，$M_t=-\log(|t|^2/r^4)$ とする．
$V$ は $v$ を正方向に回す周期 $1$ の生成場とし，
$\mu_t=(\mathcal A\lambda_t)(V)$ は同稿 \texttt{pair:cor:global-primitive}
の canonical moment とする．
定数 $m_0,c,C>0$ と $\vartheta\in(0,1/2)$ を固定できて，
$m_0\le m\le\vartheta M_t$ のすべての深度について，全 $Q$ 上で
\[
 \frac c m\le\mu_t\big|_{s=m}\le\frac C m,
 \qquad
 -\frac C m\le\mu_t\big|_{s=M_t-m}\le-\frac c m
\]
が成り立つ．従って任意の $m=m_t$ がこの範囲にあれば，
共通 interval $I_t=[-c/(2m_t),c/(2m_t)]$ の各 level は
$m_t<s<M_t-m_t$ 内で各 annulus に一つの円を持つ．
その商は $I_t\times Q$ であり，$Q$ が境界付き領域なら境界も保つ．
特に $\mu_t=0$ の level は両側深度が少なくとも $\vartheta M_t$
である領域に入る．
\end{proposition}

\begin{proof}
局所正則 frame で $\Omega=dv/v\wedge dz$ をとる．
$\bar F_t=\mathcal A\log(\omega_t^2/|\Omega|^2)$ とすると，
定数倍および base-only factor を除いた primitive の規格化から
$\mu_t=-2\pi\partial_s\bar F_t$ を得る．
Hermitian norm による $s$ と正則 frame の座標による深度の差は
固定 compact 上で有界な base-only 関数であり，$\partial_s$ は変わらない．
有限個の少し大きい積座標に命題\ref{ke:neck-integrated-normal-energy}で使った
体積包絡を適用して，$d(s)=\min(s,M_t-s)$ とすれば
ある固定 $a_0>0$ について
\[
 |\bar F_t(s,z)+2\log d(s)|\le C_0
 \qquad (a_0\le s\le M_t-a_0)
\]
を $Q$ 全体に一様に得る．Base coordinate の変更による加法項は
有限個の compact patches 上で有界であり，差商には現れない．

$2\log\Lambda>2C_0+1$ となる固定 $\Lambda>1$ を選び，
$\vartheta=1/(2\Lambda)$ とする．$2a_0\le m\le\vartheta M_t$ なら，
$m$ と $\Lambda m$ は左側半区間にある．凸性により
\[
 \partial_s\bar F_t(m,z)
 \le\frac{\bar F_t(\Lambda m,z)-\bar F_t(m,z)}{(\Lambda-1)m}
 \le -\frac{1}{(\Lambda-1)m}.
\]
逆向きの大きさの上界は命題\ref{ke:neck-integrated-normal-energy}の
差商評価 $|\partial_s\bar F_t|\le C_1/m$ による．
右側は $M_t-s$ を使って同様に求める．
$\partial^2_s\bar F_t>0$ なので $\mu_t$ は各 annulus の深度に対し
厳密に減少する．中間値定理と陰関数定理で，$I_t$ 全体にわたる
一意な深度断面が滑らかに貼り合わさる．半径を正に変えても
円束の標識は保たれるので商は $I_t\times Q$ になる．
最後は $m=\vartheta M_t$ を選ぶ．
\end{proof}

この interval は normal 深度を動かしつつ，すべての固定残余点で
同じ canonical moment を使うためのものである．
既存の固定 end の moment 漸近を moving cut に代入する必要はない．
レベル上で点ごとの円エネルギーを抑える評価と actual Moser buffer，
および基底を source や三重点へ動かす場合の一様性は別に必要である．

\begin{corollary}[軌道平均エネルギーを用いる Moser 速度評価]
\label{ke:neck-averaged-speed}
補題\ref{ke:short-circle-actual-buffer} の Moser 方程式において，
点 $x$ で
$E(x)=\int_0^1 |V|_{g_\omega}^2(\rho_a x)\,da$ とすると，
\[
 |Y_\tau(x)|_{g_\omega}
 \le \frac{\|w\|_{L^\infty}\sqrt{E(x)}}{2\sqrt{\tau(1-\tau)}} ,
 \qquad 0<\tau<1 .
\]
作用不変な開集合上で $E\le e^2$ が一様に成り立てば，
補題\ref{ke:short-circle-actual-buffer} と
系\ref{ke:short-circle-image-buffer} の長さ・buffer の結論は
$\varepsilon=e$，$c_w=(\pi/2)\|w\|_{L^\infty}$ で成り立つ．
\end{corollary}

\begin{proof}
$g_a=\rho_a^*g_\omega$ とすると
\[
 |\beta(X)|\le
 \left(\int_0^1 |w(a)|^2 g_a(V,V)\,da\right)^{1/2}
 |X|_{g_\eta}\le \|w\|_{L^\infty}\sqrt{E(x)}\,|X|_{g_\eta}.
\]
この定数を補題\ref{ke:short-circle-actual-buffer} の平方完成に用いる．
また $E\le e^2$ 上で一周の長さが $e$ 以下なので，
二点間の短い方の軌道弧の長さは $e/2$ 以下である．
これを使って両向きの軌道追跡と飽和近傍の証明をそのまま適用できる．
\end{proof}

積分評価\eqref{ke:eq:normal-energy-mass}は所与の深度での点ごとの上界や，
内側集合から小エネルギー領域の境界までの actual 距離下界を与えない．
例えば残余方向を固定しても凸関数の二階微分は小さい区間に集中し得る．
次に必要なのは，所要の開 collar を含む一様な $E$ の上界と actual buffer，
または軌道ごとに速度評価を積分して同じ被覆を保証する変位評価である．

\begin{corollary}[縮小する actual buffer からの Moser 被覆]
\label{ke:neck-shrinking-buffer}
上記の annulus--disc 座標で，compact な飽和集合 $K_t$ を含む飽和開集合
$U_t\subset\{d(s)>m_t\}$ を取る．$g_t$ に関する全曲面の距離で
\[
 \delta_t=\min\bigl(1,\operatorname{dist}_{g_t}(K_t,X_t\setminus U_t)\bigr)
\]
と置く．$m_t\delta_t^2\to\infty$ ならば，十分小さい $t$ で
$K_t$ を source と像の双方に含む飽和 source を持つ normal Moser map が存在する．
その円軌道像の ambient 直径は $O((m_t\delta_t)^{-1})$ である．
\end{corollary}

\begin{proof}
$N_t=\{\operatorname{dist}_{g_t}(\cdot,K_t)<\delta_t/2\}$ と置く．
各 $q\in N_t$ の actual $\delta_t/4$-ball の閉包は $U_t$ に含まれる．
残余射影への localized Chern--Lu 評価を半径 $\delta_t/4$ で用いると，
中心で $p^*h_D\le C\delta_t^{-2}\omega_t$ を得る．
この半径依存性は，距離 cutoff の勾配二乗と Laplacian の下界を
$O(\delta_t^{-2})$ で制御する Ricci 下界による距離比較から従う．
関連 KE 主稿の \texttt{neck:lem:rank-one-circle} と体積上界により
$N_t$ 上で $|V|_{g_t}\le\varepsilon_t=C'/(m_t\delta_t)$ となる．

$D_t=\{\operatorname{dist}_{g_t}(\cdot,K_t)<\delta_t/4\}$ とする．
$\overline D_t$ から任意の位相まで周期の半分以下で円軌道を辿るとき，
$N_t$ の最初の退出までの長さは $\varepsilon_t/2$ 以下である．
これは $\delta_t/4$ より短いので，$\overline D_t$ を通る全軌道が $N_t$ に留まる．
従って $N_t$ に含まれる最大の飽和開集合に $\overline D_t$ が含まれ，
そこで同じ短円上界を使える．$K_t$ から $\partial D_t$ までの actual 距離は
$\delta_t/4$ 以上であり，$m_t\delta_t^2\to\infty$ により
$(2c_w+1/2)\varepsilon_t$ を上回る．系\ref{ke:short-circle-image-buffer}と
補題\ref{ke:short-circle-actual-buffer}の直径評価を適用する．
\end{proof}

\subsection{固定残余領域での actual tensor 比較と Moser 被覆}

\begin{lemma}[大域 capped density を用いた射影計量の下界]
\label{ke:capped-projective-lower}
固定次数 $d\ge5$ の一般位置全平面退化を考え，
$\Omega_{t,d}^{\rm low}$ と $h_{t,d}^{\rm low}$ を
関連 KE 主稿の \texttt{coarse:cor:degree-d-volume-lower} で定める．
$\varphi_t=\log(\omega_t^2/\Omega_{t,d}^{\rm low})$ とすると，
ある固定 $A,C>0$ に対し
\[
 \operatorname{tr}_{\omega_t}(\omega_{\mathrm{FS}}|_{X_t})
 \le C\exp(A\varphi_t)
\]
が全 $X_t$ で成り立つ．従って $\varphi_t$ が一様に上から有界な領域では
$g_t\ge c g_{\mathrm{FS}}|_{X_t}$ となる．
\end{lemma}

\begin{proof}
KE 側の capped density の証明で $A_0$ を十分大きく固定して，
$h_{t,d}^{\rm low}\ge c_0\omega_{\mathrm{FS}}|_{X_t}$ を得る．
実際，背景は $(d-4)\omega_{\mathrm{FS}}|_{X_t}$ で，
すべての cap の負部分と frame errors の和は
その $O_d(A_0^{-1})$ 倍に抑えられ，残りは半正である．
同じ大域体積下界から $\varphi_t\ge-C_0$ が全 $X_t$ で成り立つ．
$T_t=\operatorname{tr}_{\omega_t}(\omega_{\mathrm{FS}}|_{X_t})$ とすると，
正則包含 $X_t\to\PP^3$ の Chern--Lu 不等式は
$\Delta_{\omega_t}\log T_t\ge-C_1-C_2 T_t$ を与える．
ここでは射影空間の固定した curvature upper bound と
$\operatorname{Ric}(\omega_t)=-\omega_t$ を用いる．
$dd^c\varphi_t=\omega_t-h_{t,d}^{\rm low}$ なので
\[
 \Delta_{\omega_t}(\log T_t-A\varphi_t)
 \ge (Ac_0-C_2)T_t-(C_1+2A).
\]
$Ac_0>C_2$ と取り，compact $X_t$ 上の最大点で $T_t$ を上から抑える．
そこでの $\varphi_t\ge-C_0$ と合わせ，
$\log T_t-A\varphi_t\le C'$ を全域に得る．
\end{proof}

\begin{theorem}[固定 source-free base の moving strip での actual Moser 被覆]
\label{ke:fixed-base-neck-moser}
関連 KE 主稿の \texttt{pair:prop:bundle-conic} の正則束座標を用い，
source と三重点を避けた固定した $Q\Subset B$ 上で
$uv=t$，$s=-\log(\|v\|^2/r^2)$，
$M_t=-\log(|t|^2/r^4)$，$d(s)=\min(s,M_t-s)$ とする．
$Q$ は compact とする．$m_t\to\infty$，$m_t\le M_t/2$ とし，
$K_t=\{p\in Q,\ d(s)\ge m_t\}$ とする．
$K_t$ を source と像の双方に含む normal Moser map が存在し，
その軌道像の actual ambient 直径は $O(m_t^{-1})$ である．
さらに $K_t$ を含む深部開 strip 上で，局所積座標の complete
annulus--disc 計量 $g_{\rm prod}$ と $g_t$ は一様に二側比較可能である．
\end{theorem}

\begin{proof}
まず $\overline Q$ の少し大きい固定残余近傍を有限個の局所積座標で覆う．
各 patch で outer annulus--disc を固定し，その内部に
三重に入れ子になった annulus--disc を，正の固定な normal radii と residual radii
で選ぶ．最も外側の入れ子領域（outer chart には compact に含まれる）では
outer chart の体積 envelope と capped density が一様に比較可能である．
これは \texttt{neck:prop:volume-envelope} の証明にある
ordinary-pair capped density の比較である．同命題の体積上界から
そこでは $\varphi_t\le C$ であり，補題\ref{ke:capped-projective-lower}
によって $g_t\ge c g_{\mathrm{FS}}$ を得る．

内側の入れ子領域から中間領域の外へ出る経路は，その最初の退出までに，
残余座標又は $u,v$ の一方を固定した正の幅だけ変化させる．
これらの座標は固定な ambient 座標上で有界な微分を持つ．
実際，正則束座標は $t$ に依存しない ambient 変換で作られ，
コンパクトな座標範囲では変換とその逆の微分が一様に有界である．
よって経路は比較が成立する部分だけで一様に正の actual 長さを持つ．
従って内側の各点の一様な半径の actual ball の閉包は中間領域に含まれる．

この actual ball 上で outer annulus--disc の complete product への
正則恒等写像に localized Chern--Lu 評価
（関連 KE 主稿 \texttt{coarse:lem:localized-chern-lu}）を適用する．
曲率 $-1$ の二曲面の積の holomorphic sectional curvature は
$-1/2$ 以下なので，中心で $g_t\ge c,g_{\rm prod}$ を得る．
既存の体積上界と合わせれば複素次元 $2$ の行列の固有値比較により
$g_t\le C g_{\rm prod}$ も従う．有限被覆で定数を一様にできる．
固定 norm による $s$ は局所積座標の深度と有界な base-only 項だけ異なり，
その項の微分も一様に有界である．従って十分深い領域上で
\[
 |V|_{g_t}\le C/d(s),\qquad
 |dp|_{g_t}\le C,\qquad
 |d\log d(s)|_{g_t}\le C
\]
となる．最後の不等式は $s=M_t/2$ の継目では片側微分，
又は経路に沿う Lipschitz 評価として読む．$|dp|$ の標的には固定基底計量を用いる．

$Q\Subset B'\Subset B$ を固定して，Moser の外側領域を
$U_t=\{p\in B',\ d(s)>m_t/2\}$ とする．
$K_t$ から $U_t$ の境界へ向かう任意の経路は，最初の退出以前に
基底で固定した正の距離を移動するか，$\log d(s)$ を $\log 2$ 以上変える．
従って actual buffer は正の定数以上である．ここで局所比較のための有限被覆は
$\overline{B'}$ 上に選んでいる．同じ議論で $B'$ と $m_t/2$ を少し内側へ
縮めた比較的 compact な領域を $U_t$ 内に取れる．
$U_t$ 上の $|V|_{g_t}=O(m_t^{-1})$ と
系\ref{ke:short-circle-image-buffer} により求める Moser map を得る．
直径は補題\ref{ke:short-circle-actual-buffer}の評価による．
\end{proof}

定理\ref{ke:fixed-base-neck-moser} は residual base が固定した場合を扱う．
Source や三重点へ近付く residual collars では，射影計量で測った入れ子領域間の幅や
capped density と complete envelopes の比較を別途調べる必要がある．
固定 compact の結論と命題\ref{ke:neck-deep-moment-strip} を合わせると，
同じ固定基底上で moving canonical moment strip を actual form に移せる．
複数 end を同時に接着するには，正則束作用と triple 作用との
開 overlap 上での整合も必要である．

\subsection{残る円の平均と同一 normal moment}

今度は $\rho_{1,s}$ を neck の領域 $U$ 全体に与えた normal 円作用とする．
End の開領域 $E\subset U$ は $\rho_1$ に飽和しており，
$E$ 上には $\rho_1$ と可換な自由な正則円作用 $\rho_2$ があり，
両者で自由な $T^2$ 作用を生成すると仮定する．
各平均と homotopy を $\mathcal A_j,K_j$ と書き，
\[
 \eta=\mathcal A_1\omega,\qquad
 \kappa=\mathcal A_2\eta=\mathcal A_1\mathcal A_2\omega,\qquad
 \lambda_1=\mathcal A_1\lambda,\qquad
 \lambda_{12}=\mathcal A_2\lambda_1
\]
と置く．$\rho_2$ は $E$ のみに必要である．

第一の Moser flow $\phi_s$ は補題\ref{ke:circle-exact-transport}を
normal action に適用して $U$ 内に構成する．
第二の flow $\psi_s$ は $E$ 内で
\begin{equation}\label{ke:eq:residual-moser-path}
 \eta_s=(1-s)\kappa+s\eta,\qquad
 \beta_2=K_2\eta=\mathcal A_1K_2\omega,\qquad
 \iota_{Z_s}\eta_s=\beta_2,\quad \psi_0=\mathrm{id}
\end{equation}
によって構成する．以下では $\psi_s$ の定義域を $T^2$ に飽和した内側領域に取り，
$\psi_1$ の像が $\phi_s$ の共通定義域に含まれる buffer を課す．
$\phi=\phi_1,\ \psi=\psi_1,\ \Theta=\phi\circ\psi$ と書く．

\begin{theorem}[共通 normal circle と全作用周期を保つ二段階平均]
\label{ke:staged-neck-averaging}
上記の仮定の下で，$\psi_s$ は $\rho_1$ と可換であり，
\[
 \Theta^*\omega=\kappa,\qquad
 \Theta^*\lambda=\lambda_{12}+dq
\]
となる一価関数 $q$ がある．元の各 $T^2$ 軌道 $\mathcal O$ に対して
$\Theta(\mathcal O)$ は actual $\omega$-Lagrangian torus であり，
neck 全体の同じ actual 円作用
$\widehat\rho_1=\phi\rho_1\phi^{-1}$ の軌道で飽和している．

Normal moment を
$\mu=\lambda_1(V_1)$，
$\widehat\mu=\mu\circ\phi^{-1}$ と規格化すると，
\begin{equation}\label{ke:eq:staged-moment-identity}
 \mu\circ\psi=\mathcal A_2\mu,\qquad
 \widehat\mu\circ\Theta=\mathcal A_2\mu.
\end{equation}
より一般に，整数 cycle $b=(b_1,b_2)\in\ZZ^2$，
$V_b=b_1V_1+b_2V_2$ に対し，
元の軌道上のその一周を $\gamma_{b,x}$ とすると
\begin{equation}\label{ke:eq:staged-period-identity}
 \int_{\Theta(\gamma_{b,x})}\lambda
 =\int_{\gamma_{b,x}}\lambda_{12}
 =\lambda_{12}(V_b)(x).
\end{equation}
従って各 end で action の加法定数を独立に取り直す必要はない．
\end{theorem}
\begin{proof}
$\eta,\kappa,\beta_2$ は $\rho_1$ に不変である．
非退化形式に関する vector field の一意性から $Z_s$ は $\rho_1$-invariant であり，
flow の一意性により $\psi_s$ は $\rho_1$ と可換である．
補題\ref{ke:circle-exact-transport}を各段に用いると
\[
 \phi^*\omega=\eta,\quad \psi^*\eta=\kappa,\qquad
 \phi^*\lambda=\lambda_1+dq_1,\quad
 \psi^*\lambda_1=\lambda_{12}+dq_2.
\]
したがって $q=q_2+q_1\circ\psi$ とすれば所要の exact identity となる．
$\lambda_{12}$ は全 $T^2$ に不変なので，可換な生成場に対して
\[
 \kappa(V_1,V_2)
 =V_1\lambda_{12}(V_2)-V_2\lambda_{12}(V_1)
     -\lambda_{12}([V_1,V_2])=0.
\]
実次元二の自由な軌道は $\kappa$-Lagrangian であり，
$\Theta^*\omega=\kappa$ によりその像も Lagrangian である．
さらに
\[
 \widehat\rho_{1,s}(\Theta x)
 =\phi\rho_{1,s}\psi x
 =\phi\psi\rho_{1,s}x
 =\Theta\rho_{1,s}x
\]
なので，同じ actual normal circle が各 torus に含まれる．

Moment の等式は加法定数まで直接計算する．
$\iota_{V_1}\eta=-d\mu$ と可換性から
\[
 \beta_2(V_1)=K_2(d\mu)=\mathcal A_2\mu-\mu.
\]
$\mu_s=(1-s)\mathcal A_2\mu+s\mu$ と置けば
$\iota_{V_1}\eta_s=-d\mu_s$ である．従って
\[
 \frac d{ds}(\mu_s\circ\psi_s)
 =\bigl(\mu-\mathcal A_2\mu+d\mu_s(Z_s)\bigr)\circ\psi_s
 =\bigl(\mu-\mathcal A_2\mu+\beta_2(V_1)\bigr)\circ\psi_s=0.
\]
初期値は $\mathcal A_2\mu$ なので
\eqref{ke:eq:staged-moment-identity}を得る．
最後に exact differential の閉 cycle 上の積分は零であり，
$\lambda_{12}(V_b)$ は軌道上で一定なので
\eqref{ke:eq:staged-period-identity}が従う．
\end{proof}

\begin{corollary}[同じ normal action を持つ開 overlap 上の一致]
\label{ke:staged-normal-overlap}
Neck と triple chart に同じ $\omega$ と同じ normal 円作用 $\rho_1$ を与え，
同じ homotopy kernel を用いるとする．
内側 overlap からの normal Moser trajectories と，
residual stage の像からの trajectories が
$\rho_1$ に飽和した共通の開 overlap に留まると仮定する．
このとき両側の normal flow $\phi_s$ と actual normal action は一致する．
Triple 側に定理\ref{ke:staged-neck-averaging}の tori を用いれば，
それらは neck 側の同じ actual normal circle を含む．
また primitive $\lambda$ と cycle markings も同じものを使えば，
moment と作用周期は\eqref{ke:eq:staged-moment-identity}，
\eqref{ke:eq:staged-period-identity}によって共通に規格化される．
独立に固定した neck foliation の第二の周期や fiber 自体の一致には，
さらに reduced leaves を同じものに選ぶ必要がある．
\end{corollary}
\begin{proof}
共通 overlap では $\mathcal A_1\omega$ と $K_1\omega$ が一致する．
従って normal Moser vector field も一致し，
trajectory がその overlap に留まるという仮定の下で ODE の一意性を使える．
残りは定理\ref{ke:staged-neck-averaging}による．
\end{proof}

\begin{lemma}[短い作用に対する二段階構成の変位]
\label{ke:staged-displacement}
Normal action に不変な Hermitian metric $h$ を $U$ に取り，
$E$ 上では全 $T^2$ に不変とする．
$ch\le g_\omega\le Ch$，
$|V_1|_h\le e_1$ on $U$，
$|V_2|_h\le e_2$ on $E$ と仮定する．
上の二段階 flow を収める内在的 $h$-buffer があるなら，
\begin{equation}\label{ke:eq:staged-displacement}
 d_h(x,\Theta x)\le\frac{C}{2c}(e_1+e_2).
\end{equation}
各元の torus が内側領域に含まれる場合，その像の ambient diameter は
\begin{equation}\label{ke:eq:staged-ambient-diameter}
 \operatorname{diam}_{g_\omega}\Theta(\mathcal O)
 \le\sqrt C(e_1+e_2)\left(\frac Cc+\frac12\right).
\end{equation}
Centered homotopies を用いれば各段の変位定数 $C/(2c)$ は $C/(4c)$ に改善できる．
\end{lemma}
\begin{proof}
平均と凸結合は同じ二側比較を満たす．
$|K_1\omega|_h\le Ce_1/2$，
$|K_2\eta|_h\le Ce_2/2$ なので，各段の vector field の長さは
$Ce_j/(2c)$ 以下である．二つの trajectory を順に結べば
\eqref{ke:eq:staged-displacement}となる．
元の同じ torus の二点は，各 circle を高々半周ずつ進むことで
$h$-長さ $(e_1+e_2)/2$ 以下の経路で結べる．
その経路の両端に Moser trajectories を付け，
$g_\omega\le Ch$ を用いると\eqref{ke:eq:staged-ambient-diameter}を得る．
Centered kernel の絶対値の積分は $1/4$ なので最後の主張も従う．
\end{proof}

\subsection{一段階平均との Hamiltonian 比較}

二段階の endpoint map は，一段階の Moser map と一般に同一ではない．
ただし exact primitive を共通に固定した場合には，次の比較ができる．

\begin{proposition}[平均の順序と整数基底に関する Hamiltonian 比較]
\label{ke:staged-hamiltonian-comparison}
$E$ 上の同じ $T^2$ 作用と $\omega=d\lambda$ を用いる．
一段階の homotopy $H=K_1+\mathcal A_1K_2$ による map を $\Phi$，
上記の二段階 map を $\Theta$ とする．
以下の path homotopy の全 flow と inverse compositions を収める共通 buffer があるなら，
共通のより小さい開領域上で $\Phi^{-1}\Theta$ は
$\kappa$ に関する Hamiltonian isotopy の endpoint である．
従って $\Theta\Phi^{-1}$ は $\omega$ に関する Hamiltonian map である．

同じ結論は，$T^2$ の任意の二つの固定した整数基底で作る一段階または
二段階の平均の間にも成立する．
Centered homotopies を用い，$\omega$ が全 $T^2$ に不変な飽和 collar がある場合，
比較 isotopy はその collar を点ごとに固定できる．
\end{proposition}
\begin{proof}
三角形 $0\le v\le u\le1$ 上で
\[
 \Omega_{u,v}=(1-u)\kappa+(u-v)\eta+v\omega,\qquad
 \Lambda_{u,v}=(1-u)\lambda_{12}+(u-v)\lambda_1+v\lambda
\]
を定める．係数が非負で和が $1$ なので $\Omega_{u,v}$ は K\"ahler であり，
$d\Lambda_{u,v}=\Omega_{u,v}$ である．
$(0,0)$ から $(1,1)$ への対角 path は一段階構成，
$(0,0)\to(1,0)\to(1,1)$ の二辺は二段階構成に対応する．
二辺は各区間の端で stationary になるよう滑らかに parameter 化する．
この path と対角 path の凸結合によって，両座標が単調な滑らかな path family
$(u_r(s),v_r(s))$ を得る．

$\beta_1=K_1\omega$，$\beta_2=K_2\eta$ とし，
各 path 上の Moser primitive を
\[
 B_{r,s}=\dot u_r(s)\beta_2+\dot v_r(s)\beta_1
\]
とする．Homotopy identity より
\begin{align*}
 B_{r,s}
 &=-\partial_s\Lambda_{u_r,v_r}
       -d\bigl(\dot u_r K_2\lambda_1+\dot v_rK_1\lambda\bigr),\\
 dB_{r,s}&=-\partial_s\Omega_{u_r,v_r}.
\end{align*}
補題\ref{ke:circle-exact-transport}と同じ微分計算により，
各 endpoint map $F_r$ は
\begin{equation}\label{ke:eq:staged-path-exactness}
 F_r^*\omega=\kappa,\qquad F_r^*\lambda=\lambda_{12}+dQ_r
\end{equation}
を満たし，$Q_r$ は $r$ に滑らかである．
両座標の総変化は各々 $1$ なので，補題\ref{ke:staged-displacement}の
総変位の上界を全 $r$ に共通に使える．
さらに $F_r(D_{\mathrm{small}})\subset F_0(D_{\mathrm{large}})$ を保証する
image buffer を取ることが，相対写像の定義に必要である．

$W_r=(\partial_rF_r)\circ F_r^{-1}$ と置く．
\eqref{ke:eq:staged-path-exactness}を $r$ で微分すると
\[
 F_r^*(\iota_{W_r}\omega)
 =d\bigl(\partial_rQ_r-F_r^*(\lambda(W_r))\bigr).
\]
従ってこの deformation は Hamiltonian であり，
$F_0^{-1}F_r$ は $\kappa$-Hamiltonian isotopy となる．
端点は $F_0=\Phi$，$F_1=\Theta$ である．

整数基底を変えても全 Haar average $\mathcal A=\mathcal A_1\mathcal A_2$ は同じである．
二つの homotopies $H,H'$ について
\[
 H\omega=\mathcal A\lambda-\lambda-d(H\lambda),\qquad
 H'\omega=\mathcal A\lambda-\lambda-d(H'\lambda)
\]
なので，その差は exact である．
同じ線形 path $(1-s)\kappa+s\omega$ に対して primitive を
$(1-r)H\omega+rH'\omega$ と補間すれば，
再び endpoint の exact identity と上の Hamiltonian 計算が成立する．
二段階の場合は既に証明した比較を前後に組み合わせればよい．

最後に，centered homotopy は invariant form を零に写す．
指定 collar 上では $\beta_1=\beta_2=H\omega=H'\omega=0$ なので，
すべての vector fields と endpoint maps はそこで零および恒等となる．
従って比較 isotopy も同 collar を点ごとに固定する．
\end{proof}

\subsection{指定した normal subgroup を持つ actual closed corner}

\begin{corollary}[二段階 closed-corner bundle と同じ GH 誤差]
\label{ke:staged-closed-corner-gh}
系\ref{ke:actual-closed-corner}の固定次数 $d\ge5$ の族と exact triple chart を用いる．
その標準 $T^2$ の任意の固定した primitive normal subgroup $\alpha$ を選び，
$\alpha,\beta$ が $\ZZ^2$ の整数基底になるよう $\beta$ を補う．
$L=-\log|t|^2$，$R=\log L$ とし，
第一段を $\alpha$，第二段を $\beta$ とする上記の二段階 map $\Theta_t$ を取る．
十分小さい $t\ne0$ についてこの map は既存の softmax domain
$\mathcal H_t=\mathcal U_t(H_{L,R})$ の開近傍に定義できる．
\[
 \mathcal W_{t,J}^{\mathrm{st}}=\Theta_t(\mathcal H_t),\qquad
 \pi_{t,J}^{\mathrm{st}}=P_t\circ\Theta_t^{-1}
\]
とすると，その全 fibers は actual Lagrangian tori であり，
同じ actual normal action
$\widehat\rho_\alpha=\phi_t^\alpha\rho_\alpha(\phi_t^\alpha)^{-1}$ を含む．

模型 $\gamma_t=g_L^{\mathrm{sf}}$ に対し
\begin{equation}\label{ke:eq:staged-actual-metric}
 (1-C/R)\gamma_t\le\Theta_t^*g_t\le(1+C/R)\gamma_t.
\end{equation}
$R^{-2}g_t$ による $\mathcal W_{t,J}^{\mathrm{st}}$ の内在距離，
および $\mathcal V_t=\{\min\ell_a>R/2\}$ 内の全経路を許す距離のどちらについても
\begin{equation}\label{ke:eq:staged-closed-distortion}
 \sup_{x,y\in\mathcal W_{t,J}^{\mathrm{st}}}
 \left|d_t(x,y)-d_J(\pi_{t,J}^{\mathrm{st}}x,\pi_{t,J}^{\mathrm{st}}y)\right|
 \le C\frac{\log R}{R}.
\end{equation}
Fiber 自身の非正規化 intrinsic diameter は $O(R^{-1})$，
正規化後は $O(R^{-2})$ である．
Raw map と retraction に関する
\eqref{ke:eq:actual-closed-corner-target}と同じ誤差評価も成立する．
定数には固定した整数基底への依存を許す．

この bundle は前節の bundle と同一の写像としては主張しない．
両者は三重点の共通の開近傍上で actual $\omega_t$-Hamiltonian map により比較でき，
その map は同じ softmax base label を持つ fibers を対応させる．
\end{corollary}
\begin{proof}
補題\ref{ke:negative-exact}により三重点の actual form は exact である．
$h=h_t$ を constrained cusp metric とする．
外側領域 $\min\ell_a>R/2$ での全深部比較と模型の一様幾何から
\[
 ch\le\gamma_t\le Ch,\qquad
 |\nabla^h\gamma_t|_h\le C,\qquad |g_t-\gamma_t|_h\le C/R,
\]
\[
 |\omega_t|_h+|\nabla^h\omega_t|_h\le C,\qquad
 |V_\alpha|_h+|\nabla^hV_\alpha|_h
 +|V_\beta|_h+|\nabla^hV_\beta|_h
 \le C/R.
\]
生成場の評価は\eqref{ke:eq:full-corner-generator-c1}の固定整数線形結合である．
$h$ は全 torus に不変なので，$\eta_t,\kappa_t$ と両 convex paths も
$h$ に関する同じ二側比較および一様 $C^1$ bounds を持つ．
Isometry の pullback は $\nabla^h$ と可換であり，
chain contraction の微分には
$|\nabla^hV|_h|\omega|_h+|V|_h|\nabla^h\omega|_h$ だけが現れる．
従って各段の primitive とその一階微分は $O(R^{-1})$ であり，
非退化な Moser 方程式とその共変微分から
\[
 |Y_{t,s}|_h+|\nabla^hY_{t,s}|_h
 \le C/R
\]
を得る．これは補題\ref{ke:open-cell-moser-c1}および
\eqref{ke:eq:full-corner-moser-c1}と同じ計算である．

Constrained cusp norm では各 $d\log\ell_a$ が一様有界である．
従って $\min\ell_a>R$ からの二つの trajectories は，
$R$ が大きければ $\min\ell_a>3R/4$ に留まり，
$\min\ell_a>R/2$ との間の固定幅の buffer により全時刻で存在する．
必要な inverse compositions と path homotopy にも，少し大きい内側領域を用いて
同じ buffer を与えられる．

各段で，$\nabla^h\gamma_t$ の上界と上記の $C^1$ 評価により
$|\mathcal L_Y\gamma_t|_{\gamma_t}\le C/R$ であるから，
Gronwall の評価と二つの flow の合成により
$\Theta_t^*\gamma_t=(1+O(R^{-1}))\gamma_t$ を得る．
Actual/model 比較を最後に引き戻せば
\eqref{ke:eq:staged-actual-metric}となる．
模型の全閉 corner の距離評価
\eqref{ke:eq:closed-corner-model-input}と一様な正規化 diameter を使えば，
内在距離について\eqref{ke:eq:staged-closed-distortion}が従う．

保護領域内の距離には，系\ref{ke:closed-corner-protected-gh}と同じ
clipped projection $\Psi$ を用いる．二つの trajectories の模型長さの和は $C/R$ 以下であり，
\eqref{ke:eq:protected-pruning-length}から
\[
 d_J(\Psi(\Theta_tz),\Psi(z))\le C/R^2
\]
となる．保護領域全体の actual/model 下比較と，
既に得た内在距離の上界を合わせれば同じ distortion bound となる．
Fiber intrinsic diameter は，模型の phase torus の直径 $C/R$ と
\eqref{ke:eq:staged-actual-metric}から従う．
Raw map の評価は補題\ref{ke:raw-values}と
補題\ref{ke:raw-lipschitz}を二つの短い trajectories に用いればよい．

最後に，一段階・二段階および整数基底の変更に対して
命題\ref{ke:staged-hamiltonian-comparison}を適用する．
共通のソース形式は全 Haar average $\mathcal A\omega_t$ であり，
同じ softmax torus の二つの像を対応させる actual Hamiltonian map を得る．
\end{proof}

\subsection{Reduced foliations による延長と残る幾何学的入力}

\begin{corollary}[共通 moment strip 上の条件付き neck 延長]
\label{ke:staged-reduced-extension}
定理\ref{ke:staged-neck-averaging}の normal 平均後の
$(U,\eta,\rho_1,\mu)$ において，共通の compact moment interval $I$ を取り，
reduced surfaces が境界付き曲面 $P$ の滑らかな束として同定されているとする．
さらに $P$ に，境界 coorientations と内部の embedded marking を固定した
基準 marked Morse foliation $\mathcal F_0$ があり，
境界近傍では radial，内部では simple saddles のみを持ち，
一つの critical leaf に一つの saddle だけを含むと仮定する．
各境界成分の互いに素な開 collar に，二段階構成から得た reduced circle foliations があり，
基準の radial circles をそれらへ写す collar 同定を $I$ に滑らかに選べるとする．
この同定は基準の collar 標識と coorientations に適合し，
互いに素な少し広い collars に台を持つ $\Diff_0(P)$ の滑らかな族へ延長できると仮定する．
このとき，それらの開 collar germs と一致する marked simple-saddle foliation family を
$P$ 上へ延ばせる．その円持ち上げを $\phi$ で移すと，
actual $\omega$-Lagrangian neck fibration が得られ，end では上記の staged tori と一致する．
その regular fiber は $T^2$，simple-saddle fiber は $S^1\times\figeight$ である．
\end{corollary}
\begin{proof}
仮定した $\Diff_0(P)$ の族により $\mathcal F_0$ を移す．
得られる族は指定した開 end germs と一致し，
内部の marked decomposition と一 critical-leaf 一 saddle の性質を保つ．
補題\ref{neck:marked-interpolation}と同様に，parameter germs の追加の補間が必要なら
collar 標識を固定した isotopies を用いる．
各固定 $\mu$ の reduced surface では一次元葉は isotropic であり，
その normal circle lift は Lagrangian である
（補題\ref{neck:lift}の reduction の証明）．
最後に $\phi^*\omega=\eta$ を用いて actual form へ移す．
\end{proof}

\begin{remark}[Actual neck への適用条件]
\label{ke:staged-neck-scope}
正則 normal action と K\"ahler 性から，
moment は各 complexified normal orbit に沿って狭義単調になる．
しかし，それだけでは全 base points に共通な moment interval や，
退化 parameter に一様な flow buffer は得られない．
系\ref{ke:staged-reduced-extension}では，これらの geometric identifications を仮定している．
本研究の actual KE 族では，source を除いた固定 compact base 上の正則 normal model に対し，
単一の global canonical primitive と共通 moment interval は
関連 KE 主稿\cite[\texttt{pair:cor:global-primitive}，
\texttt{pair:cor:actual-common-moment}]{KELocalCurrent}で検証される．
前者は base の nowhere-vanishing holomorphic $1$-form $\tau$ に対する
$\Omega=d\log v\wedge\tau$ と
$\lambda_t=d^c\log(\omega_t^2/(\Omega\wedge\overline\Omega))$ を用い，
$(\mathcal A_1\lambda_t)(V_1)$ が canonical normal moment と一致する．
後者は固定した end radius の後で $t\to0$ とする評価である．
従って，この actual 適用で残る条件は neck 全体の normal Moser flow の生存と buffer，
triple の各 end における同一作用・reduced foliation の同時一致である．
Critical actions の直線性を保つには，さらに節\ref{ng8:section}の
relative area と marking の条件が必要である．

系\ref{ke:staged-closed-corner-gh}は，三重点の標準 torus 内の任意の
指定 primitive subgroup に適合できることを示す．
別々に構成した neck action が，指定した開 collar でその subgroup と一致すること，
一つの三重点に入る複数の neck でこの一致を同時に満たすこと，
および source・branch・core の既存大域構成へ接着できることは含めない．
非定数 holomorphic unit による orbit partition の変更も，
単なる整数基底変更の Hamiltonian 比較では処理されない．
全曲面の距離と singular fibers の一様な計量評価は引き続き別の入力である．
\end{remark}
